\begin{center}
\textcolor{red}{\textbf{\LARGE Ném phần này chỗ khác}}
\end{center}
\section{Phân tích hiệu năng vận hành lý thuyết}

\subsection{Thông lượng dữ liệu (Data Throughput)}

\subsubsection{Thông lượng truyền thông nội bộ Inter-MCU}
Kênh giao tiếp giữa Master (WAN-MCU) và Slave (LAN-MCU) sử dụng SPI, xung nhịp đồng bộ lấy theo giới hạn của thiết bị Slave là $f_{\text{clk}} = 60\,\text{MHz}$.
\begin{itemize}
    \item \textbf{Thông lượng lý thuyết ($T_{\text{raw}}$):}
    $T_{\text{raw}} = f_{\text{clk}} \times 1\,\text{bit/clock} = 60\,\text{MHz} \times 1 = 60\,\text{Mbps}$
    (tương đương $7{,}5\,\text{MB/s}$).

    \item \textbf{Thông lượng hiệu dụng theo giao thức ($T_{\text{protocol}}$):}
    Sau khi tính overhead khung tin (Header, Sequence Number, CRC), cơ chế bắt tay GPIO Handshake, trễ ngắt ISR và chu trình ACK/NACK, hiệu suất giao thức đạt $\eta \approx 95\%$:
    \[
        T_{\text{protocol}} = \eta \times T_{\text{raw}} = 0{,}95 \times 60\,\text{Mbps} = 57\,\text{Mbps}.
    \]

    \item \textbf{Thông lượng duy trì thực tế ($T_{\text{eff}}$):}
    Hệ thống hoạt động theo chu kỳ batch: LAN-MCU tích lũy dữ liệu trong $t_{\text{batch}} = 10\,\text{ms}$ rồi mới thực hiện một phiên truyền SPI. Với batch payload điển hình $L = 2048\,\text{bytes} = 16384\,\text{bits}$:
    \[
        t_{\text{tx}} = \frac{L}{T_{\text{protocol}}} = \frac{16384}{57 \times 10^{6}} \approx 287\,\mu\text{s},
    \]
    \[
        T_{\text{cycle}} = t_{\text{batch}} + t_{\text{tx}} = 10\,\text{ms} + 0{,}287\,\text{ms} \approx 10{,}287\,\text{ms},
    \]
    \[
        T_{\text{eff}} = \frac{L}{T_{\text{cycle}}} = \frac{16384\,\text{bits}}{10{,}287 \times 10^{-3}\,\text{s}} \approx 1{,}59\,\text{Mbps}.
    \]
    Đây là thông lượng duy trì (sustained throughput) của kênh Inter-MCU trong điều kiện vận hành thực tế. Do lưu lượng dữ liệu cảm biến từ LAN-MCU thường nhỏ hơn $1\,\text{Mbps}$, kênh SPI không trở thành điểm nghẽn cổ chai của hệ thống.
\end{itemize}

\subsubsection{Thông lượng lưu trữ cục bộ (SDIO 4-bit)}
Giao tiếp với thẻ nhớ Micro SD hỗ trợ chuẩn SDIO 4-bit, tần số hoạt động tối đa $80 \, \text{MHz}$.
\begin{itemize}
    \item \textbf{Thông lượng lý thuyết:} $T_{\text{SDIO}} = 80 \, \text{MHz} \times 4 \, \text{bits} = 320 \, \text{Mbps}$ (tương đương $40 \, \text{MB/s}$).
    \item \textbf{Hiệu suất thực tế:} Thực nghiệm ghi nhận tốc độ ghi đạt chuẩn Class 10 (tối thiểu $10 \, \text{MB/s} \approx 80 \, \text{Mbps}$). Tốc độ này nhanh hơn gấp 8 lần thông lượng dữ liệu thực tế từ các module LAN, đảm bảo hệ thống không bị nghẽn khi thực hiện lưu đệm (Store-and-forward) lúc mất mạng WAN.
\end{itemize}
\subsubsection{Thông lượng kết nối WAN}
Phiên bản 1.0.0 hỗ trợ ba giao diện WAN với thông lượng lý thuyết khác nhau:
\begin{itemize}
    \item \textbf{Wi-Fi (802.11n, 2.4\,GHz):} Thông lượng lý thuyết tối đa $150 \, \text{Mbps}$. Trong môi trường thực tế (nhiễu RF, khoảng cách), thông lượng hiệu dụng thường đạt $20\text{--}50 \, \text{Mbps}$.
    \item \textbf{Ethernet (W5500, SPI):} Chuẩn 10/100BASE-T với thông lượng tối đa $100 \, \text{Mbps}$. Đây là giao diện ổn định nhất, phù hợp với môi trường công nghiệp yêu cầu độ trễ thấp và ổn định cao.
    \item \textbf{LTE 4G (SIM7600G-H):} Thông lượng lý thuyết Cat-4 là $150 \, \text{Mbps}$ downlink\,/\,$50 \, \text{Mbps}$ uplink. Trong điều kiện thực tế tại Việt Nam, thông lượng uplink (chiều truyền dữ liệu telemetry lên Cloud) thường đạt $5\text{--}20 \, \text{Mbps}$.
\end{itemize}
Với lưu lượng dữ liệu telemetry điển hình từ các module cảm biến ($< 1 \, \text{Mbps}$), cả ba giao diện đều đáp ứng dư yêu cầu băng thông. Cơ chế Internet Monitor đảm bảo chuyển đổi dự phòng tự động khi một giao diện gặp sự cố.
\subsection{Độ trễ hệ thống (End-to-End Latency)}
Độ trễ được xác định từ thời điểm cảm biến gửi dữ liệu đến khi dữ liệu được đóng gói MQTT gửi lên Cloud:
\begin{itemize}
    \item \textbf{Trễ xử lý tại LAN-MCU:} Bao gồm thời gian thu thập dữ liệu thô, kiểm tra Checksum và chuẩn hóa khung tin nội bộ.
    \item \textbf{Trễ truyền dẫn SPI:} Thời gian gói tin nằm trên bus SPI, phụ thuộc vào kích thước Payload và thông lượng thực tế $\approx 57\,\text{Mbps}$.
    \item \textbf{Trễ mạng diện rộng (WAN):} Phụ thuộc vào chất lượng kết nối Wi-Fi/Ethernet/LTE và độ phản hồi của MQTT Broker. Hệ thống sử dụng hàng đợi (Publish Queue) trong 8MB PSRAM để cách ly trễ mạng với trễ xử lý nội bộ.
\end{itemize}

\subsection{Khả năng chịu tải và Giới hạn kết nối}
Dựa trên tài nguyên phần cứng ESP32-S3 (16MB Flash, 8MB PSRAM), hệ thống đáp ứng:
\begin{itemize}
    \item \textbf{Giới hạn miền Sensing:} Hỗ trợ thu thập đồng thời từ 02 khe cắm LAN với đa dạng giao thức: Zigbee (IEEE 802.15.4), BLE Native/GATT, LoRa TDMA, RS485/Modbus, 1-Wire, I2C, SPI. Hệ thống tự động phát hiện (auto module detection) và nạp driver phù hợp từ NVS khi khởi động.
    \item \textbf{Dung lượng hàng đợi:} Với 8MB PSRAM, Gateway có khả năng lưu trữ tạm thời hàng ngàn gói tin telemetry (kích thước trung bình $1-2 \, \text{KB}$) trước khi phải ghi vào thẻ nhớ SD, đảm bảo vận hành ổn định ngay cả khi Broker Cloud gặp sự cố phản hồi chậm.
\end{itemize}